\chapter{Ejemplos}
En este capítulo veremos algunos ejemplos de iteraciones de funciones racionales
para introducirnos en la rama de la Dinámica Compleja.

\section{Transformaciones de Möbius}
Comencemos con el caso más sencillo: las funciones racionales de grado $1$.
Estas funciones se conocen como transformaciones de Möbius y su dinámica puede
describirse de forma bastante explícita.

\begin{definicion}[Transformación de Möbius]
    Una transformación de Möbius es una función racional de la forma
    \[
        R(z) = \frac{az+b}{cz+d}, \quad ad-bc \ne 0
    \] 
    donde $z$, $a$, $b$, $c$, $d$ son números complejos. Además, seguiremos los
    siguientes convenios con el punto del infinito: $R(\infty)=a/c$ y
    $R(-c/d)= \infty$.
\end{definicion}

Veamos primero la transformación de Möbius
\[
    R(z) = \frac{3z-2}{2z-1}.
\]

Esta función tiene un único punto fijo, $\zeta = 1$, y este es indiferente
porque $R'(1) = 1$. La pregunta que nos hacemos es: si iteramos en un punto
arbitrario del plano complejo, $z$, ¿cómo se comportará la secuencia $R^n(z)$?

Debido a que las transformaciones de Möbius son sencillas, se pueden encontrar
expresiones para $R^n(z)$ de manera sistemática. En este caso, se puede
comprobar que
\[
    R^n(z) = \frac{(2n+1)z-2n}{2nz-(2n-1)} =  1 + \frac{z-1}{2nz-(2n-1)}.
\]
Esta expresión permite ver con claridad qué ocurre: para $z \neq 1$, el
denominador crece linealmente con $n$, mientras que el numerador permanece
constante. Por ello,
\[
    R^n(z) \to \zeta,
\]
cuando $n \to \infty$. Es decir, en este ejemplo todas las órbitas convergen al
único punto fijo.

Consideremos ahora un ejemplo con un comportamiento distinto:
\[
    R(z) = \frac{(1+i)z + (1-i)}{(1-i)z + (1+i)}.
\]

En este caso tenemos dos puntos fijos, $\zeta_1 = 1$ y $\zeta_2 = -1$, ambos
indiferentes. Se puede comprobar que
\begin{equation}\label{eq:mobiusEj2Conj}
    R(z) = g^{-1}Sg(z),
\end{equation}
con
\begin{align*}
    S(z) &= -iz, \\
    g(z) &= \frac{1+z}{1-z}.
\end{align*}
La ecuación~\ref{eq:mobiusEj2Conj} muestra que $R$ es conjugada a la rotación
$S(z)=-iz$. Por tanto, la dinámica de $R$ se reduce a estudiar el comportamiento
de $S$: como $(-i)^4=1$, se tiene $S^4(z)=z$ para todo $z$. En consecuencia,
\[
    R^4(z)=\Id(z)=z,
\]
para todo $z \in \ExtC$. Así, todas las órbitas son periódicas de período que
divide a $4$; en particular, salvo los puntos fijos, no convergen.

Estos dos ejemplos ilustran bien la variedad de comportamientos que ya aparecen
en grado $1$: convergencia hacia un punto fijo en el primer caso y dinámica
periódica en el segundo. En general, las transformaciones de Möbius pueden
clasificarse según tengan uno o dos puntos fijos.

\begin{teorema}
Sea $R$ una transformación de Möbius, si $R$ tiene
\begin{enumerate}
    \item un punto fijo, $\zeta$, entonces para todo $z \in \ExtC$,
        $R^n(z) \to \zeta$ cuando $n$ tiende a $\infty$;
    \item dos puntos fijos, entonces para todo $z \in \ExtC$ se tiene que cuando
        $n$ tiende a $\infty$
    \begin{enumerate}
        \item $R^n(z)$ converge a uno de los puntos fijos,
        \item $R^n(z)$ se mueve cíclicamente en un conjunto finito de puntos,
        \item o $R^n(z)$ forma un subconjunto denso de un círculo o recta.
    \end{enumerate}
\end{enumerate}
\end{teorema}

\begin{proof}
    Véase \cite[4--5]{Beardon1991}.
\end{proof}

\section{Polinomios de grado 2}
% Hemos visto cómo las iteraciones de las transformaciones de Möbius están muy
% restringidas y tienen un buen comportamiento. Sin embargo, para el caso general
% de las funciones racionales esto no es cierto. Analicemos el comportamiento de
% algunos polinomios de grado 2, que a primera vista, parecería que también
% deberían tener un comportamiento amigable.

\ldots

% \begin{figure}[htbp]
%     \centering
%     % If your python figsize matches the page width, you don't even need scale options!
%     % But setting width=\textwidth is a safe fallback.
%     \includegraphics[width=0.8\textwidth]{images/newton_basins.pdf}
%     \caption{The Fatou basins of attraction for the rational function $R(z) = z^3 - 1$. The super-attracting fixed points are denoted by $\times$. Note that the Julia set forms the boundary of these basins.}
%     \label{fig:newton_basins}
% \end{figure}

% As we can see in Figure \ref{fig:newton_basins}, the boundary exhibits complex fractal geometry...